\subsection{Protocole expérimental}

\lipsum[9]

\subsubsection{Jeux de données}

\lipsum[10]

\subsubsection{Hyperparamètres}

Les hyperparamètres utilisés sont récapitulés dans la Table~\ref{tab:hyperparams}.

\begin{table}[H]
  \centering
  \caption{Hyperparamètres de la configuration expérimentale.}
  \label{tab:hyperparams}
  \begin{tabular}{lcc}
    \toprule
    \textbf{Paramètre} & \textbf{Valeur} & \textbf{Plage testée} \\
    \midrule
    Taux d'apprentissage & $10^{-3}$ & $[10^{-4}, 10^{-2}]$ \\
    Taille du lot & 64         & $\{32, 64, 128\}$ \\
    Époques          & 100        & — \\
    Régularisation $\lambda$ & $10^{-4}$ & $[10^{-5}, 10^{-3}]$ \\
    \bottomrule
  \end{tabular}
\end{table}

\subsection{Résultats}

\lipsum[11]

% Exemple de tableau de résultats comparatifs
\begin{table}[H]
  \centering
  \caption{Comparaison des performances (\%) sur le jeu de test.}
  \label{tab:resultats}
  \begin{tabular}{lccc}
    \toprule
    \textbf{Méthode} & \textbf{Précision} & \textbf{Rappel} & \textbf{F1} \\
    \midrule
    Méthode A \cite{exemple2024} & 82.3 & 79.1 & 80.7 \\
    Méthode B & 85.6 & 83.2 & 84.4 \\
    \textbf{Notre approche} & \textbf{89.1} & \textbf{87.5} & \textbf{88.3} \\
    \bottomrule
  \end{tabular}
\end{table}

\subsection{Analyse et discussion}

\lipsum[12]

\begin{remark}
  Notons que les résultats présentés dans la Table~\ref{tab:resultats} ont été obtenus avec une validation croisée à 5 plis afin de garantir la robustesse statistique.
\end{remark}

% 